\subsection{無形資産の識別と特徴付け}

無形資産は、組織の見えない側にある価値である。従業員が持つ業務知識、暗黙知、文書化された手順、標準、テンプレート、経験、文化などが含まれる。これらは、氷山の水面下の部分のように見えにくいが、ソフトウェアの適合性に大きく影響する。

\subsubsection{プロセスを識別し、事業目標を定義する}

まず、組織の業務プロセスと事業目標を理解する。文書化されたプロセスがある場合はそれを使い、ない場合は調査する。事業目標には、成長と継続、財務目標、従業員への責任、社会への責任、市場地位の維持などが含まれる。

\subsubsection{事業目標に結び付く無形資産を識別する}

次に、事業目標を支える無形資産を洗い出す。方針、業務プロセス、チェックリスト、教訓、テンプレート、標準、手順、計画、研修資料などが候補になる。無形資産ごとに重要度を付けると、どの資産が事業目標へ強く貢献するかを把握しやすい。

\subsubsection{無形資産を支えるソフトウェア製品を識別する}

どのソフトウェア製品が、どの無形資産を支えるかを対応付ける。ひとつの製品が複数の無形資産を支えることもあり、ひとつの無形資産が複数の製品によって支えられることもある。

\subsubsection{指標を定義し測定する}

無形資産の状態を判断するために、指標を定義する。品質指標は、無形資産そのものの質を測る。影響指標は、無形資産が業務プロセスや事業目標へどれだけ貢献するかを測る。指標は比較可能にするため、正規化や標準化が必要である。

\subsubsection{無形資産を特徴付ける}

無形資産は、品質と影響の二つの観点で特徴付けられる。品質も影響も良い資産は安定している。品質は良いが影響が低い資産は置換可能かもしれない。品質が低く影響も低い資産は警告状態である。品質が低いが影響が高い資産は改善しながら進化させる必要がある。

\par\smallskip
\begin{center}
\centering
\IfFileExists{assets/generated_figures/ch15/fig-ch15-06.png}{\includegraphics[width=0.96\linewidth]{assets/generated_figures/ch15/fig-ch15-06.png}}{\fbox{\parbox[c][48mm][c]{0.9\linewidth}{\centering 画像生成待ち\\無形資産の品質・影響マトリクス}}}
\par\smallskip
\refstepcounter{figure}\label{fig:ch15-06}図 \thefigure: 無形資産の品質・影響マトリクス
\end{center}
図 \thefigure は、無形資産の品質・影響マトリクスを視覚的に整理したものである。
\par\smallskip

\subsubsection{無形資産をビジネスモデルと結び付ける}

無形資産の状態をビジネスモデル上に配置すると、提案されたソフトウェアがどの資産を改善し、どの事業目標へ貢献するかを説明しやすくなる。これは、意思決定者に対して「なぜこのソフトウェア案が事業価値を生むのか」を示すために重要である。

\subsubsection{意思決定}

最終的には、開発すべきソフトウェア製品を優先順位付けし、選ぶ。判断基準には、無形資産の重要度、事業目標への影響、現在の特徴付け、競合への影響、ビジネスモデルへの影響、実装費用、実装期間、複雑さなどがある。これは典型的な複数属性意思決定である。
